\chapter{Compromissos de Pedersen e provas de intervalo}
\label{chapter:pedersen-commitments}

Bob já sabe reconhecer a chave da saída que Alice lhe destinou.
Como ficará a saber que recebeu 7? E como poderá a rede incluir
esses mesmos 7 numa soma sem os ler? A derivação partilhada do
capítulo anterior responde à primeira pergunta, cifrando o montante
para o destinatário. Para a segunda precisamos de um ponto que
represente o montante e se possa somar aos pontos das outras saídas.
O compromisso de Pedersen tem precisamente esta propriedade
\cite{Pedersen1992,MRL-0005-ringct}.

Os escalares do compromisso pertencem a um corpo finito. Uma moeda,
pelo contrário, tem um montante inteiro entre zero e um máximo:
nas saídas de RingCT estudadas, \(2^{64}-1\). Ao chegarmos à soma
dos compromissos, teremos de ligar cada escalar a um inteiro desse
intervalo. É daí que virá a necessidade de provar alguma coisa
acerca dos algarismos do montante \cite{monero-source-v16}.

\section{Compromissos}
\label{sec:commitments}
\index{compromisso@compromisso}
\index{ocultacao@ocultação}
\index{vinculacao@vinculação}

Um esquema de compromisso permite a Alice fixar uma mensagem sem a revelar de
imediato. Mais tarde, Alice apresenta uma \emph{abertura}; Bob verifica que é
a mensagem originalmente fixada. Queremos duas propriedades:

\begin{description}
 \item[Ocultação.] Antes da abertura, o compromisso não deve revelar a
 mensagem.
 \item[Vinculação.] Depois de publicado o compromisso, Alice não deve
 conseguir abri-lo como duas mensagens diferentes.
\end{description}

Para comprometer uma mensagem \(m\), Alice gera um sal aleatório longo \(s\) e
publica
\[
 h=\mathcal H\bigl(s\mathbin\|m\bigr).
\]
Na abertura, revela \((s,m)\). Bob repete o hash e compara-o com \(h\). O sal
impede a enumeração antecipada quando o domínio de mensagens é pequeno. Neste
compromisso por hash, a ocultação depende do sal e a vinculação depende, em
particular, de ser impraticável encontrar colisões. Esta construção não
preserva a adição das mensagens comprometidas.

\section{Compromissos de Pedersen}
\label{pedersen_section}
\index{compromisso@compromisso!de pedersen@de Pedersen}
\index{pedersen, compromisso de@Pedersen, compromisso de}
\index{homomorfismo dos compromissos@homomorfismo dos compromissos}

Para somar montantes ocultos, empregamos um compromisso de Pedersen
\cite{Pedersen1992}. Sejam \(G\) e \(H\) pontos não nulos do subgrupo de ordem
prima \(l\). O grupo é cíclico, portanto existe matematicamente um escalar
\(\gamma\) tal que
\[
 H=\gamma G.
\]
A expressão «geradores independentes» designa aqui a hipótese computacional
de ser impraticável obter \(\gamma\) a partir dos pontos públicos.
O gerador \(H\) usado por Monero tem a receita pública da
secção~\ref{subsec:pedersen-fixed-generator}. Essa receita permite a todos
reconstituir o mesmo ponto; a dificuldade de recuperar a relação entre os
dois geradores é a hipótese que sustenta a vinculação.

Para uma máscara \(x\in\mathbb Z_l\) e um montante escalar
\(b\in\mathbb Z_l\), definimos
\begin{equation}
 C(x,b)=xG+bH.
 \label{eq:pedersen-commitment}
\end{equation}
O par \((x,b)\) é uma abertura do ponto \(C\). Os escalares canónicos usados
pela implementação ocupam 32 bytes em ordem \emph{little-endian} e representam
inteiros menores que \(l\). Um montante Monero é um inteiro sem sinal de
64 bits, colocado nos oito bytes menos significativos de um escalar; os outros
24 bytes são nulos. Os pontos têm a codificação comprimida de Edwards de
32 bytes.


\textbf{Ocultação.} Se \(x\) for uniforme em \(\mathbb Z_l\), então, para
qualquer \(b\), o ponto \(xG+bH\) é uniforme no subgrupo. O esquema abstracto
tem assim ocultação perfeita: mesmo cálculo ilimitado não escolhe uma abertura
preferida entre todas as possibilidades \cite{maxwell-ct,
SCOZZAFAVA1993313}. Na carteira, as máscaras de saída são derivadas de um
segredo Diffie--Hellman por um hash reduzido; para um observador, a sua
pseudo-aleatoriedade é uma garantia computacional, não uma fonte literal de
aleatoriedade perfeita.

\textbf{Vinculação.} Se Alice encontrasse duas aberturas distintas
\((x,b)\ne(x',b')\) para o mesmo ponto, obteria
\[
 (x-x')G=(b'-b)H.
\]
Salvo o caso trivial em que ambas as diferenças são nulas, isso revelaria a
relação entre \(G\) e \(H\). A vinculação é portanto \emph{computacional}, sob
a dificuldade do logaritmo discreto. Quem
conhecesse \(\gamma\) poderia escolher qualquer \(b'\) e calcular
\(x'=x+(b-b')\gamma\), quebrando-a.

\textbf{Homomorfismo aditivo.} A aritmética dos escalares é módulo~\(l\), e
\begin{align}
 C(x,a)+C(y,b)&=C(x+y,a+b),\label{eq:commitment-addition}\\
 C(x,a)-C(y,b)&=C(x-y,a-b).\label{eq:commitment-subtraction}
\end{align}
Isto deixa comparar somas ocultas. Não significa que a igualdade pública de
dois compromissos revele a igualdade dos montantes: máscaras diferentes podem
compensar montantes diferentes. Para provar uma relação entre montantes é
preciso também ligar ou cancelar as máscaras de maneira verificável.

\subsection{A construção do gerador de montantes}
\label{subsec:pedersen-fixed-generator}
\index{h@\(H\), gerador de montantes}
\index{gerador@gerador!de montantes@de montantes}

O ponto fixo \(H\) obtém-se a partir da codificação convencional de \(G\).
Se \(\operatorname{dec}\) designar a descompressão de Edwards da
secção~\ref{point_compression_section}, a receita é
\begin{equation}
 Q=\operatorname{dec}\!\left(
           \operatorname{Keccak256}(\operatorname{enc}(G))\right),
 \qquad H=[8]Q.
 \label{eq:pedersen-fixed-generator}
\end{equation}
Neste argumento concreto, os 32 bytes do resumo constituem uma codificação
válida. O último byte é \texttt{e1}: o seu bit superior dá a paridade
ímpar de \(x_Q\), e os restantes sete bits entram na ordenada. A sequência
que representa \(y_Q\) termina, portanto, em \texttt{61}.
Obtém-se \(x_Q\) pela raiz de
\[
 x_Q^2=\frac{y_Q^2-1}{dy_Q^2+1}
\]
com a paridade indicada. A descompressão determina assim \(Q\) sem qualquer
escolha secreta.

Três duplicações, efectuadas pela lei de adição da curva, dão os pontos
seguintes. Cada linha apresenta a codificação comprimida do ponto à esquerda:
\begin{center}
\small
\begin{tabular}{r@{\quad}l}
\(Q\)&\texttt{1e609dc09abbd7bca35029ffbf164d05496b441a97d9f13020d869762acd0ee1}\\
\([2]Q\)&\texttt{5e80078173851a4aeb1563aa6e0da45729dc453e7ca922bab939d011e57af412}\\
\([4]Q\)&\texttt{1cdf7620c9ba31e2f191f689d1bfbec2ee5d111c126a289d8ded73da456f82a6}\\
\([8]Q\)&\texttt{8b655970153799af2aeadc9ff1add0ea6c7251d54154cfa92c173a0dd39c1f94}
\end{tabular}
\end{center}
A última linha fixa \(H\) \cite{monero-point-maps}. Temos
\([l]H=[8l]Q=\mathcal O\); a codificação distingue-o da identidade.
Como \(l\) é primo, a ordem de \(H\) é exactamente \(l\).

Esta construção e a função \(\mathcal H_p^{(1)}\) da
Equação~\eqref{eq:hash-to-point-one} usam o mesmo resumo de \(G\) em
operações diferentes. Aqui os bytes codificam directamente \(Q\).
Na função hash para pontos, todos os bits são lidos como um inteiro,
reduzidos módulo \(q\) e entregues ao mapa \(M\). Por isso, a base
\(\mathcal H_p^{(1)}(G)\), cuja codificação começa por \texttt{d632},
e o gerador \(H\), cuja codificação começa por \texttt{8b65}, são
pontos distintos. A validade desta descompressão concreta basta para
definir o gerador fixo; o mapa \(M\) fornece a operação definida para
todos os argumentos de \(\mathcal H_p^{(1)}\).

\section{Montantes ocultos}
\label{sec:pedersen_monero}
\index{montante oculto@montante oculto}
\index{mascara@máscara!de um compromisso@de um compromisso}

Consideremos a saída \(t\), de montante \(b_t\), enviada a Bob. A Secção
\ref{sec:one-time-addresses} definiu a derivação partilhada \(\mathcal D_B\)
e o escalar dependente do índice
\[
 h_t=\mathcal H_n\bigl(\mathcal D_B\mathbin\|
              \operatorname{varint}(t)\bigr).
\]
Para um endereço principal ordinário, \(\mathcal D_B=8rK_B^v\) do lado de
Alice e \(\mathcal D_B=8k_B^vR\) do lado de Bob. Sub-endereços, troco e chaves
adicionais escolhem a chave pública de transacção apropriada como descrito no
capítulo anterior; depois de obtido \(h_t\), o tratamento do montante é o
mesmo.

A carteira calcula a máscara de compromisso
\begin{equation}
 y_t=\mathcal H_n\bigl(
       \texttt{commitment\_mask}\mathbin\|h_t\bigr)
 \label{eq:commitment-mask}
\end{equation}
e publica o compromisso normal
\begin{equation}
 C_t^b=C(y_t,b_t)=y_tG+b_tH.
 \label{eq:output-commitment}
\end{equation}
Aqui \(y_t\) é um escalar, não um ponto. A redução do hash produz a sua
codificação módulo~\(l\); a construção não rejeita explicitamente a máscara
nula, cuja ocorrência tem probabilidade desprezável no modelo do hash.

O montante propriamente dito cabe em oito bytes. Seja
\[
 q_t=\operatorname{primeiros8}\!\left(
   \mathcal H(\texttt{amount}\mathbin\|h_t)\right).
\]
Alice coloca no campo \texttt{ecdhInfo.amount}
\begin{equation}
 e_t=\operatorname{LE8}(b_t)\mathbin\oplus q_t,
 \label{eq:encoded-amount}
\end{equation}
onde \(\oplus\) é o XOR da Secção \ref{sec:XOR_section}. Desde o formato
\texttt{RCTTypeBulletproof2}, incluindo CLSAG e Bulletproofs+, só estes oito
bytes são serializados; a máscara não é transmitida porque ambos os lados a
recompõem pela
Equação~\eqref{eq:commitment-mask}. O formato anterior, mais volumoso, guardava
32 bytes de montante e uma máscara codificada; a versão~10 introduziu a forma
compacta sob o nome histórico \texttt{RCTTypeBulletproof2}
\cite{monero-source-v16}.


Bob calcula o mesmo \(h_t\), recupera
\[
 \operatorname{LE8}(b_t)=e_t\mathbin\oplus q_t
\]
e recompõe \(y_t\). A carteira exige que as representações escalares obtidas
sejam canónicas, que os 24 bytes superiores do montante sejam nulos e que
\[
 y_tG+b_tH=C_t^b.
\]
Uma chave de vista permite assim detectar a saída, descodificar o montante e
confirmar a abertura; não permite gastá-la. O XOR não oferece segurança de
teoria da informação nem autenticação autónoma: a confidencialidade depende
computacionalmente do segredo partilhado e do hash, enquanto o compromisso
público fornece a verificação de consistência da abertura.


\section{Introdução a RingCT}
\label{sec:ringct-introduction}
\index{ringct@RingCT}
\index{pseudo-compromisso@pseudo-compromisso}

Uma saída reúne objectos com funções diferentes. A chave de saída de
utilização única liga a saída à autorização de gasto; \(C_t^b\) oculta o
montante; e \(e_t\) permite à carteira certa recuperar esse montante. Nenhum
destes objectos, isoladamente, oferece anonimato completo nem prova que a
transacção está autorizada.

Suponhamos que a transacção gasta \(m\) saídas anteriores de montantes
\(a_1,\ldots,a_m\), cria \(p\) saídas de montantes
\(b_0,\ldots,b_{p-1}\) e declara a taxa pública \(f\). A conservação desejada
é
\begin{equation}
 \sum_{j=1}^{m}a_j=\sum_{t=0}^{p-1}b_t+f.
 \label{eq:amount-balance}
\end{equation}

O verdadeiro compromisso gasto pela entrada \(j\) tem a forma
\[
 C_j^a=x_jG+a_jH.
\]
Publicá-lo como a entrada verdadeira denunciaria o membro escolhido no anel.
Por isso, o remetente cria um \emph{pseudo-compromisso de saída}
\[
 \widetilde C_j^a=x'_jG+a_jH.
\]
Para o membro verdadeiro,
\begin{equation}
 C_j^a-\widetilde C_j^a=(x_j-x'_j)G.
 \label{eq:pseudo-output-difference}
\end{equation}
O Capítulo \ref{chapter:transactions} mostrará como a CLSAG prova, sem revelar
o membro, conhecimento da chave de gasto e desta diferença de máscaras. Um
pseudo-compromisso não é, só por ser publicado, uma prova de propriedade ou
de igualdade de montantes.

Os compromissos das novas saídas são os da
Equação~\eqref{eq:output-commitment}. Para os primeiros \(m-1\)
pseudo-compromissos, a implementação escolhe máscaras independentes e
uniformes em \(\mathbb Z_l^*\). A última é determinada por
\begin{equation}
 x'_m=\sum_{t=0}^{p-1}y_t-
       \sum_{j=1}^{m-1}x'_j\pmod l.
 \label{eq:last-pseudo-mask}
\end{equation}
Logo, \(\sum_jx'_j=\sum_ty_t\), e o verificador testa a equação de pontos
\begin{equation}
 \boxed{\displaystyle
 \sum_{j=1}^{m}\widetilde C_j^a
   =\sum_{t=0}^{p-1}C_t^b+fH.}
 \label{eq:commitment-balance}
\end{equation}


Ao expandi-la, o componente em \(G\) cancela por construção e sobra a
Equação~\eqref{eq:amount-balance}, interpretada módulo~\(l\). O público vê os
compromissos e \(f\), mas não as suas aberturas. Esta verificação prova o
balanço dos compromissos; a CLSAG trata da autorização e da escolha anónima da
entrada, e as provas de intervalo que se seguem excluem montantes de saída
modulares ilegítimos.

\section{Provas de intervalo}
\label{sec:range_proofs}
\index{prova de intervalo@prova de intervalo}
\index{prova de dominio@prova de domínio|see{prova de intervalo}}
\index{intervalo de 64 bits@intervalo de 64 bits}

A Equação~\eqref{eq:commitment-balance} é uma igualdade no subgrupo de
ordem~\(l\). Sem outra condição, ela aceita escalares que não são montantes
Monero. Por exemplo, tomemos entradas de 6 e 5 unidades e saídas formalmente
iguais a 21 e \(-10\). No corpo escalar,
\[
 (6+5)-\bigl(21+(l-10)\bigr)=-l\equiv0\pmod l.
\]
Se as máscaras também se cancelassem, a equação de compromissos passaria,
apesar de \(l-10\) não ser um montante admissível. A redução modular torna
a igualdade compatível com representantes inteiros que dariam a aparência
de criar valor.

Uma prova de intervalo exclui esses representantes modulares. Para cada compromisso de saída
\(C_t^b\), a afirmação exacta é
\begin{equation}
 \exists\,y_t\in\mathbb Z_l,\quad
 \exists\,b_t\in\{0,\ldots,2^{64}-1\}:
 \quad C_t^b=y_tG+b_tH.
 \label{eq:range-statement}
\end{equation}
O montante \(b_t\) e a máscara \(y_t\) são as testemunhas privadas; o
compromisso, os geradores e a prova são públicos. Os extremos são inclusivos:
há exactamente \(2^{64}\) valores possíveis, todos representáveis nos oito
bytes usados pela carteira.

Esta prova não identifica o destinatário, não prova propriedade, não autoriza
o gasto, não impede uma dupla despesa e não verifica por si só o balanço da
transacção. Prova apenas a relação de intervalo para os compromissos que entram
no seu traslado. A validação RingCT combina-a depois com a equação de balanço e
com as assinaturas das entradas.

\section{Assinaturas em anel Borromean}
\label{sec:borromean}
\index{borromean, assinatura@Borromean, assinatura}
\index{assinatura em anel@assinatura em anel!borromean@Borromean}

As primeiras transacções RingCT provaram a
Equação~\eqref{eq:range-statement} com assinaturas em anel Borromean
\cite{Signatures2015BorromeanRS,MRL-0005-ringct}. A construção permite ver o
intervalo nascer, bit a bit, de compromissos e assinaturas de Schnorr. É um
formato histórico: a versão~9 deixou de o admitir em novas transacções, embora
continue a ser necessário verificar as provas antigas da lista de blocos.


Seja o compromisso ao montante \(b\):
\[
 C_a=xG+bH.
\]
Primeiro, fazemos uma partição aleatória da máscara \(x\), tal que
\[
 x=\sum_{i=0}^{63}x_i\pmod l.
\]
Podemos escolher os primeiros 63 escalares e determinar o último pela
diferença. Podemos também escolher todas as máscaras e definir \(x\) pela
soma; em ambos os casos, é esta igualdade que liga as partes ao compromisso.

Segundo, descrevemos \(b\) em termos binários:
\[
 b=\sum_{i=0}^{63}b_i2^i,\qquad b_i\in\{0,1\}.
\]
Para cada bit, formamos um compromisso
\begin{equation}
 C_i=x_iG+b_i2^iH.
 \label{eq:borromean-bit-commitment}
\end{equation}
O compromisso inicial fica então decomposto:
\begin{equation}
 \begin{split}
 \sum_{i=0}^{63}C_i
   &=\left(\sum_{i=0}^{63}x_i\right)G
     +\left(\sum_{i=0}^{63}b_i2^i\right)H\\
   &=xG+bH=C_a.
 \end{split}
 \label{eq:borromean-sum}
\end{equation}

O objectivo é provar que o montante de \(C_a\) pertence ao intervalo. Mas
\(C_a\), só por si, não prova tal coisa: é somente um ponto de curva
elíptica. Publicar os \(C_i\) e verificar a soma também não basta. Falta
provar que cada compromisso contém uma das duas parcelas permitidas,
\(0\) ou \(2^i\), sem revelar qual delas.

Se \(b_i=0\), temos
\[
 C_i=x_iG+0H=x_iG.
\]
Se \(b_i=1\), temos
\[
 C_i-2^iH=x_iG+2^iH-2^iH=x_iG.
\]
Como tal, fixamos para cada índice um anel de duas chaves públicas:
\begin{equation}
 P_{i,0}=C_i,\qquad P_{i,1}=C_i-2^iH.
 \label{eq:borromean-ring-pair}
\end{equation}
Alice conhece a chave privada \(x_i\) de uma delas. O índice dessa chave é
precisamente o bit \(b_i\). Uma assinatura em anel prova conhecimento de uma
das duas chaves privadas e esconde o índice; a composição borromeana reúne
os 64 anéis num único desafio.

\subsection{Preparar os dois casos}

Vamos seguir a geração e depois refazer, como Bob, as contas da verificação.
Todos os escalares que se seguem pertencem a \(\mathbb Z_l\), e todas as
operações entre eles são feitas módulo~\(l\). Empregamos \(\mathcal H_n\)
para obter um escalar a partir das codificações dos pontos, como nas provas
anteriores.

Alice escolhe um escalar aleatório novo \(\alpha_i\) para cada anel. Primeiro
prepara um ponto final \(\psi_i\); só depois de preparar todos os anéis pode
calcular o desafio comum. Aqui os dois casos começam por caminhos diferentes.

\paragraph{Se \(b_i=0\).}
Alice conhece a chave privada de \(C_i\). Escolhe também uma resposta
aleatória \(\beta_{i,1}\) para a outra chave e calcula
\begin{align*}
 \omega_i&=\mathcal H_n(\alpha_iG),\\
 \psi_i&=\beta_{i,1}G+\omega_i(C_i-2^iH).
\end{align*}
O ponto \(\alpha_iG\) começa o percurso na chave que Alice conhece; o
desafio \(\omega_i\) permite avançar pela outra chave, com a resposta
simulada \(\beta_{i,1}\). A resposta da primeira chave fica por calcular.

\paragraph{Se \(b_i=1\).}
Alice conhece a chave privada de \(C_i-2^iH\). Prepara simplesmente
\[
 \psi_i=\alpha_iG.
\]
Agora o ponto final está na chave que ela conhece. A passagem pela primeira
chave e as duas respostas serão calculadas depois do desafio comum.

\subsection{O desafio que liga os 64 anéis}

Uma vez preparados todos os pontos finais, Alice calcula
\begin{equation}
 \lambda=\mathcal H_n(\psi_0\mathbin\|\psi_1\mathbin\|\cdots
                         \mathbin\|\psi_{63}).
 \label{eq:borromean-common-challenge}
\end{equation}
A ordem é a dos bits, do menos significativo ao mais significativo, e cada
ponto entra através da sua codificação de 32 bytes. O desafio \(\lambda\)
fica assim preso aos resultados de todos os anéis. Não é um escalar que Alice
possa escolher livremente depois de ver as respostas.


Falta fechar cada percurso. Se \(b_i=0\), Alice calcula a resposta que
deixara em aberto:
\[
 \beta_{i,0}=\alpha_i-x_i\lambda.
\]
Se \(b_i=1\), escolhe uma resposta aleatória \(\beta_{i,0}\), passa pela
primeira chave e fecha na segunda:
\begin{align*}
 \mu_i&=\beta_{i,0}G+\lambda C_i,\\
 \theta_i&=\mathcal H_n(\mu_i),\\
 \beta_{i,1}&=\alpha_i-x_i\theta_i.
\end{align*}

Alice envia os 64 compromissos \(C_i\) e a assinatura
\[
 \sigma_{\rm B}=
 \bigl(\lambda,(\beta_{i,0})_{i=0}^{63},
                 (\beta_{i,1})_{i=0}^{63}\bigr).
\]
Os bits \(b_i\), as máscaras \(x_i\) e os escalares \(\alpha_i\) ficam
privados. Os pontos \(\psi_i\) não são enviados: Bob irá reconstruí-los.
Também não se enviam os \(2^iH\), pois são pontos fixos que o verificador
conhece de antemão.

\subsection{A mesma verificação para zero e para um}

Bob começa por verificar a igualdade \(C_a=\sum_iC_i\). Depois, para cada
índice, faz sempre as mesmas contas:
\begin{align*}
 \mu_i'&=\beta_{i,0}G+\lambda C_i,\\
 \theta_i'&=\mathcal H_n(\mu_i'),\\
 \psi_i'&=\beta_{i,1}G+\theta_i'(C_i-2^iH).
\end{align*}
No fim, calcula
\[
 \lambda'=\mathcal H_n(\psi_0'\mathbin\|\psi_1'\mathbin\|\cdots
                          \mathbin\|\psi_{63}')
\]
e exige \(\lambda'=\lambda\). Vamos ver por que ambos os casos chegam à
mesma igualdade.

\paragraph{Quando o bit é zero.}
Sabemos que \(C_i=x_iG\) e que
\(\beta_{i,0}=\alpha_i-x_i\lambda\). Portanto,
\begin{align*}
 \mu_i'
   &=\beta_{i,0}G+\lambda C_i\\
   &=(\alpha_i-x_i\lambda)G+\lambda C_i\\
   &=\alpha_iG-\lambda(x_iG)+\lambda C_i\\
   &=\alpha_iG.
\end{align*}
Os dois termos em \(\lambda\) cancelam-se. Logo,
\[
 \theta_i'=\mathcal H_n(\alpha_iG)=\omega_i,
\]
e o ponto final reconstruído é
\begin{align*}
 \psi_i'
   &=\beta_{i,1}G+\theta_i'(C_i-2^iH)\\
   &=\beta_{i,1}G+\omega_i(C_i-2^iH)\\
   &=\psi_i.
\end{align*}

\paragraph{Quando o bit é um.}
A primeira resposta foi escolhida aleatoriamente. Bob reconstrói exactamente
o ponto que Alice usou para calcular a segunda resposta:
\[
 \mu_i'=\beta_{i,0}G+\lambda C_i=\mu_i,
 \qquad \theta_i'=\theta_i.
\]
Agora usamos \(C_i-2^iH=x_iG\) e
\(\beta_{i,1}=\alpha_i-x_i\theta_i\):
\begin{align*}
 \psi_i'
   &=\beta_{i,1}G+\theta_i'(C_i-2^iH)\\
   &=(\alpha_i-x_i\theta_i)G+\theta_i(C_i-2^iH)\\
   &=\alpha_iG-\theta_i(x_iG)+\theta_i(x_iG)\\
   &=\alpha_iG=\psi_i.
\end{align*}
É na segunda chave que os termos se cancelam. Em ambos os casos,
\(\psi_i'=\psi_i\); repetindo o argumento para os 64 índices, obtemos
\[
 \lambda'
  =\mathcal H_n(\psi_0\mathbin\|\cdots\mathbin\|\psi_{63})
  =\lambda.
\]

Bob não recebeu nenhuma indicação do caso escolhido. A diferença está na
resposta que Alice conseguiu calcular com a sua chave privada; as equações
que Bob executa são idênticas. Há também um verdadeiro fecho circular: o
desafio inicial produz os pontos finais, e o hash desses pontos tem de
reproduzir o desafio inicial.

\subsection{Por que razão isto prova o intervalo}

As contas anteriores mostram que uma prova honesta é aceite. Para excluir
uma falsificação precisamos, além disso, da solidez de conhecimento da
assinatura borromeana, no modelo do oráculo aleatório, e da dificuldade do
logaritmo discreto. Sob essas hipóteses, uma assinatura válida atesta
conhecimento de um escalar \(z_i\) e de um índice \(d_i\in\{0,1\}\) em cada
anel, tais que
\[
 P_{i,d_i}=z_iG,
 \qquad C_i=z_iG+d_i2^iH.
\]
Estas testemunhas, que Bob não conhece, definem uma abertura do compromisso
inicial. Com efeito,
\begin{align*}
 C_a
   &=\sum_{i=0}^{63}C_i\\
   &=\left(\sum_{i=0}^{63}z_i\right)G
     +\left(\sum_{i=0}^{63}d_i2^i\right)H.
\end{align*}
O montante desta abertura é um inteiro entre zero e
\[
 \sum_{i=0}^{63}2^i=2^{64}-1.
\]
A vinculação computacional do compromisso impede Alice de conhecer outra
abertura com montante diferente. É assim que a assinatura prova a
Equação~\eqref{eq:range-statement}: uma das duas chaves em cada anel determina
um bit, e a igualdade da soma prende todos os bits ao mesmo compromisso.

Não basta dizer que o algoritmo de geração aceita apenas bits. Um atacante
pode escrever outro algoritmo; é a relação verificada, juntamente com as
hipóteses criptográficas, que limita o montante. Em particular, se aceitássemos
um \(\lambda\) livre sem exigir
\eqref{eq:borromean-common-challenge}, bastaria escolher as respostas e
calcular \(\mu_i'\), \(\theta_i'\) e \(\psi_i'\) pela própria fórmula de
verificação. Todas essas contas seriam possíveis sem conhecer uma única chave
privada. O teste final é precisamente o constrangimento que falta a essa
tentativa.

O conhecimento do bit permanece oculto pelas propriedades de conhecimento
zero da composição, e a máscara oculta o montante. Isto não significa que os
montantes possíveis sejam equiprováveis na vida real: a prova não elimina
informação exterior à transacção, nem escolhe a distribuição dos pagamentos.
\footnote{A derivação matemática usa pontos do subgrupo primo. O formato
legado descodifica os compromissos como pontos da curva sem um teste autónomo
de pertença a esse subgrupo. Os desafios são Keccak--256 reduzida módulo~\(l\),
sem um prefixo textual de separação de domínio. Estas regras pertencem à prova
de intervalo histórica; a assinatura de mensagens com contexto explícito do
capítulo \ref{chapter:advanced-schnorr} é uma construção distinta.}

\subsection{Os pesos e o espaço ocupado}

É aparente que o papel de \(2^i\) poderia ser desempenhado por outros pesos
públicos. Com inteiros positivos \(k_1,\ldots,k_m\), formaríamos
\[
 C_i=x_iG+b_ik_iH,\qquad
 (P_{i,0},P_{i,1})=(C_i,C_i-k_iH),
\]
e, reunindo agora \(m\) anéis no desafio comum, provaríamos uma abertura com
montante
\[
 b=\sum_{i=1}^{m}b_ik_i,\qquad b_i\in\{0,1\}.
\]
Se quiséssemos manter todas estas somas no intervalo Monero, exigiríamos
\(\sum_i k_i\leq2^{64}-1\) como desigualdade de inteiros. Mas os pesos
determinariam um conjunto de somas admissíveis, que poderia deixar lacunas.

Por exemplo, se \(k_1=2\) e \(k_2=3\), as quatro escolhas dos bits dariam
\[
 0\cdot2+0\cdot3=0,\quad
 1\cdot2+0\cdot3=2,\quad
 0\cdot2+1\cdot3=3,\quad
 1\cdot2+1\cdot3=5.
\]
Bob saberia que o montante pertence a \(\{0,2,3,5\}\), embora não soubesse
qual deles foi escolhido. Com outros pesos, escolhas distintas poderiam
ainda produzir a mesma soma: \(2^m\) conta os vectores de bits, não garante
\(2^m\) montantes distintos. As potências de dois dão-nos exactamente o que
pretendemos: todos os inteiros do intervalo, sem lacunas e com uma
representação única. Sendo fixas, também dispensam a transmissão dos pesos.

Resta contar o preço desta clareza. Para uma única saída, transmitimos 64
pontos \(C_i\), 64 respostas \(\beta_{i,0}\), 64 respostas
\(\beta_{i,1}\) e o desafio comum \(\lambda\). Cada ponto ou escalar ocupa
32 bytes, pelo que a prova histórica ocupa
\[
 (64+64+64+1)\times32=193\times32=6176\text{ bytes}.
\]
O compromisso da saída já figura na transacção e não está contado outra vez
neste total. Cada saída precisa da sua própria prova. As Bulletproofs que se
seguem conservam a decomposição em bits, mas comprimem a prova e permitem
agregar vários montantes.

\section{Bulletproofs}
\label{sec:bullet_proofs}
\index{bulletproofs@Bulletproofs}
\index{produto interno@produto interno}
\index{argumento de produto interno@argumento de produto interno}
\index{ipa@IPA|see{argumento de produto interno}}

A carteira tem um compromisso \(V=\gamma G+vH\) e conhece a sua
abertura. Para provar que \(v\) cabe em 64 bits, a construção
Borromean associava um anel a cada bit. Uma Bulletproof começa também
pelos bits. A economia vem de os comprometer em conjunto e de provar,
numa só conta, que cada um vale zero ou um e que, somados com os
respectivos pesos, dão o montante de \(V\)~\cite{Bulletproofs_paper}.

\index{peso@peso!valor posicional dos bits@valor posicional dos bits}
Comecemos onde se vêem todos os algarismos. O número cinco, escrito
com quatro bits e começando pelas unidades, dá \((1,0,1,0)\).
A verificação tem duas partes. A primeira é
\(1\cdot1+0\cdot2+1\cdot4+0\cdot8=5\). A segunda é que cada
entrada \(b\) satisfaz \(b(b-1)=0\). Num corpo, um produto só
se anula quando um dos factores é zero. Por isso, estas quatro
igualdades obrigam cada \(b\) a pertencer a \(\{0,1\}\).
Na prova, o facto de as contas se fazerem módulo \(l\) permite falar
de montantes inteiros porque \(2^{64}<l\). A soma de 64 bits com
os seus pesos nunca chega a dar a volta ao módulo.

Fixemos agora \(n=64\). Sublinhamos os vectores e numeramos as suas
posições desde zero. As coordenadas de \(\underline a_L\) são os
bits e as de \(\underline a_R=\underline a_L-\underline1\) são os
mesmos bits diminuídos de uma unidade. O símbolo \(\circ\) indica
multiplicação coordenada a coordenada, e os parênteses angulares somam
as parcelas de um produto interno. As condições a provar escrevem-se
\[
 \langle\underline a_L,\underline2^{\,n}\rangle=v,\qquad
 \underline a_L-\underline a_R=\underline1,\qquad
 \underline a_L\circ\underline a_R=\underline0 ,
 \quad \underline2^{\,n}=(1,2,\ldots,2^{n-1}).
\]
O provador escolhe dois vectores aleatórios \(\underline s_L,\underline s_R\)
e duas máscaras \(\alpha,\rho\). Com geradores vectoriais públicos
\(\underline G,\underline H\) de relações discretas desconhecidas, envia
\begin{align*}
 A&=\langle\underline a_L,\underline G\rangle+
    \langle\underline a_R,\underline H\rangle+\alpha G,\\
 S&=\langle\underline s_L,\underline G\rangle+
    \langle\underline s_R,\underline H\rangle+\rho G.
\end{align*}
O compromisso \(A\) fixa os dois vectores da afirmação. O segundo
ponto fornece a aleatoriedade com que se poderá mais tarde responder
sobre eles conservando os bits secretos.

\index{peso@peso!dos residuos de uma prova de intervalo@dos resíduos de uma prova de intervalo}
Para combinar as posições, o traslado produz um desafio \(y\ne0\).
Se as diferenças a testar forem \(e_0,\ldots,e_{n-1}\), a soma
\(e_0+e_1y+\cdots+e_{n-1}y^{n-1}\) é a avaliação de um polinómio.
Uma diferença fixa não nula dá um polinómio não nulo, com no máximo
\(n-1\) raízes. Um desafio uniforme entre os \(l-1\) escalares não
nulos encontra uma dessas raízes com probabilidade até
\((n-1)/(l-1)\). O compromisso antecede o desafio justamente para
que as diferenças estejam fixadas por ocasião da escolha.
As potências do desafio seguinte, \(z\ne0\), separam no polinómio
os três tipos de condição. Escreva-se \(\underline y^{\,n}=(1,y,\ldots,y^{n-1})\)
e formem-se
\begin{align*}
 \underline l_0&=\underline a_L-z\underline1,\\
 \underline r_0&=\underline y^{\,n}\circ
                        (\underline a_R+z\underline1)+z^2\underline2^{\,n}.
\end{align*}
Multiplicando na posição \(i\), aparece
\((a_{L,i}-z)(a_{R,i}+z)
=a_{L,i}a_{R,i}+z(a_{L,i}-a_{R,i})-z^2=z-z^2\).
A parcela que multiplica \(2^i\) é \(z^2(a_{L,i}-z)2^i\).
Somando ambas nas \(n\) posições, obtemos
\begin{equation}
 t_0=\langle\underline l_0,\underline r_0\rangle
      =z^2v+\delta(y,z),\qquad
 \delta(y,z)=(z-z^2)\sum_{i=0}^{n-1}y^i-z^3\sum_{i=0}^{n-1}2^i .
 \label{eq:bp-t0-relation}
\end{equation}
O termo \(\delta\) contém apenas números públicos. A dedução com a
notação vectorial completa encontra-se no Apêndice~\ref{appendix:nota}.

Voltemos ao montante cinco para perceber o que fica à vista no
produto. Com quatro posições e os desafios ilustrativos \(y=3,z=2\),
temos \(\underline l_0=(-1,-2,-1,-2)\) e
\(\underline r_0=(6,11,34,59)\). O produto vale
\(-6-22-34-118=-180\). O verificador calcula por si
\(\delta=(2-4)(1+3+9+27)-8(1+2+4+8)=-200\).
A diferença é \(-180-(-200)=20=2^2\cdot5\).
Este exemplo explica por que é preciso ocultar o próprio \(t_0\).
Quem recebesse \(t_0\) poderia subtrair \(\delta\) e dividir por
\(z^2\), recuperando o montante. Os números escritos com sinal
representam os seus resíduos no corpo de escalares.

O provador oferece, em vez disso, uma avaliação mascarada. Antes de
escolher o ponto de avaliação, trabalha com a indeterminada \(X\) e
os polinómios de grau um
\[
 \underline l(X)=\underline l_0+X\underline s_L,\qquad
 \underline r(X)=\underline r_0+
                 X(\underline y^{\,n}\circ\underline s_R).
\]
Chamemos \(\underline r_1=\underline y^{\,n}\circ\underline s_R\).
Por distributividade, os coeficientes de
\(t(X)=\langle\underline l(X),\underline r(X)\rangle=t_0+t_1X+t_2X^2\)
são o \(t_0\) já conhecido do provador e
\[
 t_1=\langle\underline s_L,\underline r_0\rangle+
     \langle\underline l_0,\underline r_1\rangle,\qquad
 t_2=\langle\underline s_L,\underline r_1\rangle .
\]
O provador envia \(T_1=t_1H+\tau_1G\) e \(T_2=t_2H+\tau_2G\), com
máscaras frescas. Só então se extrai o desafio \(x\). O provador
calcula \(t(x)\) e publica três escalares
\[
 t(x),\qquad \tau_x=z^2\gamma+x\tau_1+x^2\tau_2,\qquad
 \mu=\alpha+x\rho .
\]
Os pontos públicos já contêm exactamente estas combinações de
máscaras. O termo em \(\gamma\) vem do compromisso \(V\), e os
termos em \(\tau_1,\tau_2\) vêm de \(T_1,T_2\). Nas aberturas
de \(A+xS\), as máscaras somam \(\mu\). Chegámos assim às duas
igualdades que o verificador precisa de conferir.

\subsection{Uma igualdade para o montante e outra para os vectores}
\label{sec:bullet_ver}
\index{bulletproofs@Bulletproofs!duas condicoes de verificacao@duas condições de verificação}

As duas igualdades seguintes ligam o compromisso ao valor e o compromisso
aos vectores. O verificador confere a primeira com os dados públicos; o
argumento de produto interno permitirá conferir a segunda juntamente
com \(t(x)=\langle\underline l(x),\underline r(x)\rangle\), conservando
secretos os dois vectores de 64 elementos.

De $t_0=z^2v+\delta(y,z)$ e das aberturas dos $T_i$ vem
\begin{equation}
t(x)H+\tau_xG
=z^2V+\delta(y,z)H+xT_1+x^2T_2.
\label{eq:bp-value-check}
\end{equation}
Para conferir as coordenadas, convém usar os pontos
\(H'_i=y^{-i}H_i\). A multiplicação por \(y^i\) em \(r_i(x)\)
cancela-se com a divisão por \(y^i\) em \(H'_i\). Os termos restantes
são precisamente os termos públicos acrescentados a \(\underline l_0\)
e \(\underline r_0\), e por isso
\begin{equation}
\begin{split}
&\langle\underline l(x),\underline G\rangle+
 \langle\underline r(x),\underline{H'}\rangle+\mu G\\
&\qquad=A+xS+
 \langle z\underline y^{\,n}+z^2\underline2^{\,n},\underline{H'}\rangle+
 \langle-z\underline1,\underline G\rangle.
\end{split}
\label{eq:bp-vector-check}
\end{equation}

Depois de \(t(x)\) entrar no traslado, ambos os lados extraem \(x_{ip}\) e
calculam \(U=x_{ip}H\). O verificador reconstrói o ponto público
\begin{equation}
\begin{split}
 P_k={}&A+xS-\mu G+t(x)U\\
 &+\langle-z\underline1,\underline G\rangle
 +\langle z\underline y^{\,n}+z^2\underline2^{\,n},
                                      \underline{H'}\rangle.
\end{split}
\label{eq:bp-public-ip-point}
\end{equation}
A Equação~\eqref{eq:bp-vector-check} identifica os termos vectoriais
neste ponto. Ao substituir também \(t(x)\) pelo produto interno,
obtemos a relação que o provador deve demonstrar:
\[
 P_k=\langle\underline l(x),\underline G\rangle+
     \langle\underline r(x),\underline{H'}\rangle+
     \langle\underline l(x),\underline r(x)\rangle U.
\]
As duas parcelas lineares fixam os vectores comprometidos, e a parcela
em \(U\) liga-os ao valor declarado do seu produto interno.

Ponhamos \(k=\log_2 64=6\). O argumento dobra esta relação até
restarem os escalares \(a,b\) e dois geradores finais. A identidade
acumulada das dobras é
\begin{equation}
 P_k=P_0-\sum_{j=1}^{k}
              \left(u_j^2L_j+u_j^{-2}R_j\right).
 \label{eq:bp-ip-recursion}
\end{equation}
O verificador obtém \(P_k\) dos dados públicos pela
Equação~\eqref{eq:bp-public-ip-point}; os pares \(L_j,R_j\) e os
escalares finais permitem conferir esta última igualdade. A derivação
das dobras, que veremos a seguir, explica como cada par transporta os
termos cruzados de metade das coordenadas para a metade seguinte.

Nem os \(u_j\) atravessam a rede: ambos os lados os colhem do traslado. Para
uma saída enviam-se \(2\cdot6+9=21\) objectos de 32 bytes:
\[
 t(x),\tau_x,a,b,\mu,T_1,T_2,A,S,\quad
 L_1,\ldots,L_6,\quad R_1,\ldots,R_6.
\]
São 672 bytes, antes dos pequenos prefixos de serialização, em vez dos 6176
bytes da prova Borromean de uma saída.

\subsection{Como se reduzem os dois vectores a dois escalares}
\label{sec:bullet_detalhes}
\index{argumento de produto interno@argumento de produto interno!dobras@dobras}

Imaginemos que só restam duas coordenadas de cada vector. Chamemos-lhes
\(a_0,a_1\) e \(b_0,b_1\), e aos seus geradores \(G_0,G_1\)
e \(H'_0,H'_1\). O ponto a verificar tem a abertura
\[
 P_1=a_0G_0+a_1G_1+b_0H'_0+b_1H'_1+(a_0b_0+a_1b_1)U .
\]
Para ficar com uma coordenada de cada espécie, o provador vai somar
os dois \(a_i\) com pesos recíprocos e fazer o mesmo aos \(b_i\)
na ordem inversa. Sabe já que essa operação produzirá termos
\emph{cruzados}, em que uma coordenada de índice zero encontra a de
índice um. Dispõe-nos em dois pontos e envia-os antes do desafio:
\[
 L_1=a_0G_1+b_1H'_0+a_0b_1U,\qquad
 R_1=a_1G_0+b_0H'_1+a_1b_0U .
\]
Obtido o desafio \(u_1\ne0\), o provador calcula
\(a=u_1a_0+u_1^{-1}a_1\) e \(b=u_1^{-1}b_0+u_1b_1\).
O verificador pode combinar as bases nas proporções inversas e obter
\[
 G_*=u_1^{-1}G_0+u_1G_1,\qquad H'_*=u_1H'_0+u_1^{-1}H'_1.
\]
Vejamos o que resulta de multiplicar \(a\) por \(G_*\). Aparecem
\(a_0G_0\) e \(a_1G_1\), já contidos em \(P_1\), além de
\(u_1^2a_0G_1\) e \(u_1^{-2}a_1G_0\). Estes últimos pertencem a
\(u_1^2L_1\) e \(u_1^{-2}R_1\). A mesma expansão em \(bH'_*\) e
em \(abU\) dá a identidade completa
\[
 aG_*+bH'_*+abU=P_1+u_1^2L_1+u_1^{-2}R_1 .
\]
Cada participante chega ao mesmo ponto por um caminho que lhe é
acessível. O provador calcula com as aberturas. O verificador
recebe \(a,b\), reconstrói \(G_*,H'_*\) e dispõe na prova dos
dois pontos \(L_1,R_1\).

Para \(2^k\) posições a operação é a mesma. Onde escrevemos um escalar,
tomamos agora uma metade do vector. Onde multiplicámos dois escalares,
tomamos o produto interno das duas metades. A primeira dobra de uma
prova de 64 bits transforma duas listas de comprimento 64 em duas de
comprimento 32. Seguem-se \(16,8,4,2,1\). Ordenemos as rondas pelo
comprimento que lhes dá entrada, de \(j=k\) até \(j=1\). Em cada uma
o provador envia \(L_j,R_j\), o traslado produz \(u_j\ne0\), e o
ponto passa de \(P_j\) a \(P_{j-1}=P_j+u_j^2L_j+u_j^{-2}R_j\).
A soma das rondas fornece
\begin{align}
 P_k&=P_0-\sum_{j=1}^k(u_j^2L_j+u_j^{-2}R_j),
              \label{eq:bp-p-recursion}\\
 P_0&=aG_*+bH'_*+abU. \label{eq:bp-p-final}
\end{align}

\subsection{O peso de cada gerador na igualdade final}
\label{sec:bullet_fim}
\index{peso@peso!nas dobras do produto interno@nas dobras do produto interno}
\index{argumento de produto interno@argumento de produto interno!geradores finais@geradores finais}

O verificador consegue construir directamente os geradores finais. Com
quatro geradores, a primeira dobra coloca o primeiro e o terceiro no
mesmo par, depois o segundo e o quarto. A dobra seguinte reúne os dois
resultados. Por conseguinte
\[
 G_*=
 u_1^{-1}(u_2^{-1}G_0+u_2G_2)
 +u_1(u_2^{-1}G_1+u_2G_3).
\]
Olhemos para \(G_2\). O índice \(2=10_2\) situa-o na metade direita
na primeira dobra e na esquerda na segunda. O seu peso é
\(u_2u_1^{-1}\). Para qualquer número de rondas, o bit de índice
\(j-1\) diz em que metade cai o gerador na ronda \(j\). Se
\(b(i,j)\) valer \(-1\) para o bit zero e \(+1\) para o bit um, temos
\begin{equation}
 \phi_i=\prod_{j=1}^k u_j^{b(i,j)},\qquad
 G_*=\sum_i\phi_iG_i,\qquad H'_* =\sum_i\varphi_iH'_i,\qquad
 \varphi_i=\phi_i^{-1}.
 \label{eq:bp-g-arrow}
\end{equation}
Cada factor é invertido para a família \(H'\) porque a dobra dos
seus geradores usa os expoentes contrários aos dos \(G\).

O verificador recompõe os desafios pela ordem das mensagens. No formato
clássico de Monero, seja \(h_0\) o hash escalar da concatenação dos
compromissos de saída na representação dividida por oito, e escreva-se
\(\mathcal H_n\) para o hash reduzido a escalar. No hash seguinte,
os pontos designam as respectivas codificações no formato da prova.
As dependências exactas são
\begin{align*}
 y&=\mathcal H_n(h_0\|A\|S),&
 z&=\mathcal H_n(y),\\
 x&=\mathcal H_n(z\|z\|T_1\|T_2),&
 x_{ip}&=\mathcal H_n(x\|x\|\tau_x\|\mu\|t(x)).
\end{align*}
Começando no estado \(x_{ip}\), cada par \(L_j,R_j\) produz
\(u_j=\mathcal H_n(\text{estado anterior}\|L_j\|R_j)\) e o estado
seguinte é \(u_j\). A primeira condição a testar continua a ser
\eqref{eq:bp-value-check}. Para a segunda, o verificador reúne os
termos das rondas em
\[
 \Sigma_{\rm IP}=\sum_{j=1}^k(u_j^2L_j+u_j^{-2}R_j)
\]
e testa \(P_k+\Sigma_{\rm IP}=aG_*+bH'_*+abU\).
Substituindo aqui \(P_k\) da Equação~\eqref{eq:bp-public-ip-point},
\(U=x_{ip}H\) e \(H'_i=y^{-i}H_i\), obtém exactamente a mesma
igualdade expressa nos geradores originais
\begin{equation}
\begin{split}
 \mathcal O\stackrel{?}{=}&\ A+xS-\mu G+\Sigma_{\rm IP}
                         +(t(x)-ab)x_{ip}H\\
 &+\langle-z\underline1-a\underline\phi,\underline G\rangle\\
 &+\left\langle z\underline1+\underline y^{-n}\circ
       (z^2\underline2^{\,n}-b\underline\varphi),
       \underline H\right\rangle .
\end{split}
\label{eq:bp-final}
\end{equation}
Aqui \(\underline y^{-n}=(1,y^{-1},\ldots,y^{-(n-1)})\). O termo
\(-\mu G\) retira a máscara de \(A+xS\). A soma \(\Sigma_{\rm IP}\)
actualiza o ponto em cada dobra, e \((t(x)-ab)x_{ip}H\) reúne
as contribuições inicial e final do produto interno. Sob as hipóteses
de logaritmo discreto e do modelo de oráculo aleatório para
Fiat--Shamir, as duas igualdades, as condições de formato e a
reconstituição dos desafios sustentam a afirmação de que a abertura
de \(V\) contém um montante no intervalo prometido~\cite{Bulletproofs_paper}.

\section{Bulletproofs+}
\label{sec:bulletproof_plus}
\index{bulletproofs+@Bulletproofs+}
\index{bulletproof+@Bulletproof+|see{Bulletproofs+}}
\index{produto interno@produto interno!ponderado@ponderado}

Em Bulletproofs+ o compromisso inicial continua a conter os bits de
todas as saídas. A diferença algébrica está no produto de que se prova
conhecimento. Com os índices a começar em zero, cada parcela
\(a_i b_i\) entra com o peso \(y^{i+1}\). Esse peso acompanha as
dobras até à última coordenada, e duas novas mensagens, \(A_1,B\),
permitem provar a relação entre os dois escalares finais sem os
revelar~\cite{bulletproofs-plus-paper}.

As duas formas serializadas tornam a diferença visível:
\begin{align*}
 \Pi_{\rm BP}
  &=(A,S,T_1,T_2,\tau_x,\mu,
     \boldsymbol L,\boldsymbol R,a,b,t),\\
 \Pi_+
  &=(A,A_1,B,r_1,s_1,d_1,
     \boldsymbol L,\boldsymbol R).
\end{align*}
Para a mesma dimensão, Bulletproofs+ retira três objectos de 32 bytes. Em ambas as construções uma prova pode abranger vários montantes.

A ideia algébrica continua a começar em bits. Para \(p\) montantes, ponhamos
\[
 N=64,\qquad
 M=2^{\lceil\log_2 p\rceil},\qquad
 n=MN,
\]
com \(M=1\) quando \(p=1\). Concatenam-se os bits dos \(p\) montantes num
vector \(\boldsymbol a_L\in\mathbb Z_l^n\); as \(M-p\) posições de montante
restantes são preenchidas implicitamente com zeros. Definimos
\[
 \boldsymbol a_R=\boldsymbol a_L-\boldsymbol 1.
\]
As duas relações
\begin{align}
 b_j&=\left\langle
   \boldsymbol a_L^{(j)},\boldsymbol 2^{64}\right\rangle,
   \label{eq:bp-value-bits}\\
 \boldsymbol a_L\circ\boldsymbol a_R&=\boldsymbol 0
   \label{eq:bp-bit-constraint}
\end{align}
dizem, respectivamente, que cada valor é a soma dos seus bits e que cada
coordenada é 0 ou 1. Aqui
\(\boldsymbol 2^{64}=(1,2,\ldots,2^{63})\), \(\circ\) é o produto coordenada a
coordenada e \(\boldsymbol a_L^{(j)}\) é o bloco de 64 coordenadas do montante
\(j\), para \(0\leq j<p\).

O provador compromete-se com estes vectores sem os revelar. Cada dobra publica
dois pontos, pelo que são necessárias \(\log_2(64M)\) rondas em vez de 64
compromissos por saída. A verificação continua a ter trabalho proporcional a
\(64M\): a agregação poupa sobretudo espaço, não transforma toda a verificação
num cálculo logarítmico.

\subsection{Formatos de prova e representação dos pontos}

As Bulletproofs clássicas figuram em três formatos históricos.
Em \texttt{RCTTypeBulletproof2} associam-se a um montante cifrado
de oito bytes e à máscara derivada da secção~\ref{sec:pedersen_monero}.
Em \texttt{RCTTypeCLSAG} a autorização já usa CLSAG. O primeiro
formato desta família chamava-se \texttt{RCTTypeBulletproof}.
Na especialização da versão 16, \texttt{RCTTypeBulletproofPlus}
contém a segunda estrutura, a prova Bulletproofs+.

Nas equações anteriores, os pontos representam os elementos do grupo que
entram na verificação. A serialização de Bulletproofs guarda
\(A,S,T_1,T_2,\boldsymbol L,\boldsymbol R\) multiplicados por
\(8^{-1}\) no subgrupo primo; o verificador multiplica os pontos recebidos
por 8 para recuperar os elementos das equações. Bulletproofs+ usa a mesma
convenção para \(A,A_1,B,\boldsymbol L,\boldsymbol R\). Esta transformação
fixa a correspondência entre os bytes da prova e os pontos matemáticos.

Para o mesmo \(M\), a parte formada por objectos de 32 bytes ocupa
\begin{align}
 \lvert\Pi_{\rm BP}\rvert
   &=32\bigl(9+2\log_2(64M)\bigr),\label{eq:bp-size}\\
 \lvert\Pi_{+}\rvert
   &=32\bigl(6+2\log_2(64M)\bigr).\label{eq:bpp-size}
\end{align}
Não se contam aqui os pequenos prefixos de comprimento da serialização.
Bulletproofs+ retira três objectos, ou 96 bytes: para \(M=2\), são 736 contra
640 bytes; para \(M=16\), 928 contra 832 bytes. A prova Borromean, recordemos,
ocupava 6176 bytes \emph{por saída}. Estes números são comparações de formato;
não afirmam, sem uma medição separada, que um verificador seja tantas vezes
mais rápido do que outro.


\subsection{Bulletproofs+ na versão 16}

Para não sobrecarregar a notação anterior, chamemos \(v_j\) ao montante e
\(\gamma_j\) à máscara de compromisso nesta subsecção:
\[
 C_j=\gamma_jG+v_jH.
\]
Uma Bulletproofs+ agregada prova, para todos os seus \(p\) compromissos,
\begin{equation}
 \exists\,(\gamma_j,v_j)_{j=0}^{p-1}:
 \quad C_j=\gamma_jG+v_jH
 \quad\land\quad 0\leq v_j<2^{64}.
 \label{eq:bpp-statement}
\end{equation}
As aberturas \((\gamma_j,v_j)\) são as testemunhas. Os compromissos
\((C_0,\ldots,C_{p-1})\), os geradores, o traslado e
\[
 \Pi_+=(A,A_1,B,r_1,s_1,d_1,
          \boldsymbol L,\boldsymbol R)
\]
são públicos. O vector de compromissos internos \(V\) é reconstruído a partir
dos compromissos de saída, no campo \texttt{outPk.mask}; a prova serializada
contém apenas os elementos de \(\Pi_+\).

\subsubsection*{Agregação e preenchimento}

Uma transacção RingCT canónica deste tipo contém exactamente uma Bulletproofs+,
e essa prova cobre todos os seus compromissos de saída. O formato admite
\(1\leq p\leq16\); a regra geral da versão~16 exige separadamente pelo menos
duas saídas numa transacção comum. O valor interno \(M\) é a menor potência de
dois não inferior a \(p\):
\[
\begin{array}{c|ccccc}
 p&1&2&3\text{--}4&5\text{--}8&9\text{--}16\\ \hline
 M&1&2&4&8&16\\
 \lvert\boldsymbol L\rvert=\lvert\boldsymbol R\rvert
  &6&7&8&9&10
\end{array}
\]
Os montantes fictícios usados para completar \(M\) são zeros implícitos; não
criam saídas nem compromissos serializados.

Durante a análise da transacção, o tamanho de
\(\boldsymbol L,\boldsymbol R\) tem de corresponder precisamente a esse
\(M\), o número reconstruído de \(V_j\) tem de ser \(p\), e o limite é 16.
Assim, uma prova para capacidade 16 não é uma codificação canónica de apenas
duas saídas.


\subsubsection*{As duas representações do compromisso}

A estrutura de transacção guarda o compromisso normal \(C_j\). O provador,
porém, constrói
\begin{equation}
 V_j=8^{-1}C_j=(8^{-1}\gamma_j)G+(8^{-1}v_j)H,
 \label{eq:bpp-v-eighth}
\end{equation}
onde \(8^{-1}\) é o inverso do cofactor no corpo \(\mathbb Z_l\). A construção
da prova produz estes \(V_j\); a transacção guarda \(8V_j=C_j\) no campo
\texttt{outPk.mask}. Na leitura, cada compromisso serializado é multiplicado
por \(8^{-1}\) para restaurar \(V_j\).

Os pontos \(A,A_1,B,\boldsymbol L,\boldsymbol R\) são igualmente produzidos e
serializados na representação dividida por 8. Antes da equação final, o
verificador multiplica todos esses pontos e cada \(V_j\) por 8. Isso elimina
qualquer componente de torção e deixa os pontos usados na multiplicação
multiescalar no subgrupo de ordem prima. O compromisso serializado
\(C_j=8V_j\) já tem, portanto, a representação exigida pela equação final.


\subsubsection*{Geradores e traslado}

Além de \(G\) e \(H\), a prova usa vectores públicos
\(\boldsymbol G=(G_0,\ldots,G_{n-1})\) e
\(\boldsymbol H=(H_0,\ldots,H_{n-1})\). Cada ponto é derivado dos
bytes de \(H\), da cadeia ASCII \texttt{bulletproof\_plus}, da codificação
\texttt{varint} do índice e de uma função hash-para-ponto com limpeza pelo
cofactor. Uma derivação que produzisse a identidade seria rejeitada. A
segurança pressupõe que não se conhecem relações de logaritmo discreto úteis
entre estes geradores. Na notação do artigo Bulletproofs+, os papéis dos
geradores de valor e de máscara aparecem trocados: na notação presente, \(H\)
leva o montante e \(G\) leva a máscara.

O traslado começa num ponto fixo derivado da cadeia
\texttt{bulletproof\_plus\_transcript}. Escrevendo
\(\mathcal H_n\) para Keccak--256 seguido de redução módulo \(l\), a ordem é:
\begin{align}
 T_1&=\mathcal H_n\!\left(
   T_0\mathbin\|\mathcal H_n(V_0\mathbin\|\cdots\mathbin\|V_{p-1})
   \right),\notag\\
 y&=\mathcal H_n(T_1\mathbin\|A),&
 z&=\mathcal H_n(y),\label{eq:bpp-transcript-start}\\
 T^{(0)}&=z,&
 u_k=T^{(k+1)}&=\mathcal H_n\!\left(
   T^{(k)}\mathbin\|L_k\mathbin\|R_k\right),
   \label{eq:bpp-transcript-round}\\
 e&=\mathcal H_n\!\left(
   T^{(\log_2 n)}\mathbin\|A_1\mathbin\|B\right).
 \label{eq:bpp-transcript-end}
\end{align}
A cadeia inicial separa este protocolo de outros usos do hash. Todos os
desafios são escalares reduzidos; o provador recomeça e o verificador rejeita
se \(y,z,u_k\) ou \(e\) forem zero.


\subsubsection{A abertura que entra nas dobras}
\label{sec:bpp_detalhes}

Definamos o compromisso vectorial
\[
 \left\langle\boldsymbol a,\boldsymbol G\right\rangle+
 \left\langle\boldsymbol b,\boldsymbol H\right\rangle
 =\sum_{i=0}^{n-1}a_iG_i+\sum_{i=0}^{n-1}b_iH_i.
\]
O primeiro ponto enviado é
\begin{equation}
 A=8^{-1}\!\left(
   \langle\boldsymbol a_L,\boldsymbol G\rangle+
   \langle\boldsymbol a_R,\boldsymbol H\rangle+\alpha G
   \right),
 \label{eq:bpp-a}
\end{equation}
com \(\alpha\in_R\mathbb Z_l^*\). Depois dos desafios \(y,z\), o provador
constrói a janela
\begin{equation}
 d_{jN+i}=z^{\,2(j+1)}2^i,
 \quad 0\leq j<M,\quad0\leq i<N,
 \label{eq:bpp-window}
\end{equation}
e os vectores
\begin{align}
 \boldsymbol a'&=\boldsymbol a_L-z\boldsymbol1,\notag\\
 b'_i&=a_{R,i}+z+d_i\,y^{\,n-i}.
 \label{eq:bpp-prime-vectors}
\end{align}
A máscara acumulada que liga estes vectores aos compromissos é
\begin{equation}
 \alpha'=\alpha+y^{\,n+1}
       \sum_{j=0}^{p-1}z^{\,2(j+1)}\gamma_j.
 \label{eq:bpp-alpha-prime}
\end{equation}
Podemos ver na mesma parcela os três factores que o nó há-de
reconstituir. A posição \(k=jN+i\) entra no produto com \(y^{k+1}\).
O termo da janela, no vector direito, entra com \(y^{n-k}\).
Multiplicados, dão \(y^{n+1}\) para \emph{todas} as posições.
O factor que ainda distingue as saídas é \(z^{2(j+1)}\), e o que
ainda distingue os bits de cada saída é \(2^i\). Assim,
\[
 \sum_{i=0}^{N-1} a_{L,jN+i} \,y^{n+1}d_{jN+i}
   = y^{n+1}z^{2(j+1)}\sum_{i=0}^{N-1} 2^i a_{L,jN+i}
   = y^{n+1}z^{2(j+1)} b_j.
\]
É por isso que a máscara \(\gamma_j\) do compromisso de \(b_j\)
comparece em \(\alpha'\) com o mesmo coeficiente. Para as saídas
de 7 e de 3, por exemplo, os dois termos do montante dão
\(y^{129}(7z^2+3z^4)\) se a saída de Bob for a primeira. Estes
produtos fazem-se em \(\mathbb Z_l\). As potências de dois serviram
para interpretar os inteiros antes de os pesos dos desafios reunirem
as suas relações.
\index{peso@peso!no produto ponderado de Bulletproofs+@no produto ponderado de Bulletproofs+}

O argumento de produto interno dobra
\((\boldsymbol a',\boldsymbol b',
  \boldsymbol G,\boldsymbol H)\) em cada desafio \(u_k\). Cada ronda envia
\(L_k,R_k\), inclui duas máscaras aleatórias novas e reduz sucessivamente a
dimensão até 1. Chamemos \(a^\star,b^\star\) aos escalares finais,
\(G^\star,H^\star\) aos geradores finais e \(\alpha^\star\) à máscara
acumulada depois de todas as dobras.
Com escalares aleatórios não nulos \(r,s,d,\eta\), o fecho é
\begin{align}
 A_1={}&8^{-1}\!\left(
  rG^\star+sH^\star+dG+
  y(rb^\star+sa^\star)H\right),\label{eq:bpp-a1}\\
 B={}&8^{-1}\!\left(\eta G+rysH\right),\label{eq:bpp-b}\\
 r_1={}&r+ea^\star,\qquad
 s_1=s+eb^\star,\label{eq:bpp-rs}\\
 d_1={}&\eta+ed+e^2\alpha^\star.
 \label{eq:bpp-d1}
\end{align}
Na criação pela carteira, as máscaras do compromisso são derivadas como na
Equação~\eqref{eq:commitment-mask}; \(\alpha\), as máscaras das rondas e
\(r,s,d,\eta\) vêm do gerador de escalares uniformes não nulos. Esta
aleatoriedade é necessária para ocultar os bits e as aberturas. A ocultação de
Pedersen perfeita, a propriedade de conhecimento da prova e a heurística
Fiat--Shamir são afirmações distintas; a solidez e o conhecimento de
Bulletproofs+ são garantias computacionais sob as hipóteses do protocolo, não
uma consequência da ocultação perfeita por si só.

\subsection{Verificação}
\label{sec:bpp_ver}
\index{bulletproofs+@Bulletproofs+!verificacao@verificação}

O verificador não recebe os bits nem as máscaras. Recompõe \(V_j\) e todos os
desafios do traslado das Equações~\eqref{eq:bpp-transcript-start}--
\eqref{eq:bpp-transcript-end}. Recompõe também as potências de \(y,z,u_k\) e os
geradores dobrados. As respostas \(r_1,s_1,d_1\) substituem os escalares finais nas
Equações~\eqref{eq:bpp-a1}--\eqref{eq:bpp-d1}. Depois de reunir termos, toda a
prova se reduz a uma multiplicação multiescalar que tem de dar a identidade
\(\mathcal O\).

O verificador pode reunir várias provas desta família numa só equação.
Se \(\mathcal F_q\) é a combinação de pontos que deveria ser
\(\mathcal O\) para a prova \(q\), escolhe pesos independentes
\(w_q\in_R\mathbb Z_l^*\) e testa
\begin{equation}
 \sum_q w_q\mathcal F_q=\mathcal O.
 \label{eq:bpp-batch}
\end{equation}
Os pesos impedem, excepto com probabilidade desprezável da ordem de \(1/l\),
que erros de provas diferentes se cancelem de propósito. Isto é
\emph{verificação em lote} de várias provas; não é a agregação de vários
montantes dentro de uma prova.

Cada lote reúne provas da mesma família: Bulletproofs+ e Bulletproofs têm
equações próprias. O conjunto de transacções de um bloco pode fornecer um
lote; a validação isolada de uma transacção usa o mesmo procedimento com
uma única prova. A escolha do lote é independente da agregação dos montantes
que cada prova contém.


\subsubsection{Condições de aceitação}
\label{sec:bpp_fim}

Antes da equação final, há condições de aceitação que não devem desaparecer
atrás da álgebra:

\begin{itemize}
 \item \(r_1,s_1,d_1\) têm de ser codificações escalares canónicas, isto é,
 menores que \(l\). Os desafios produzidos por hash são reduzidos módulo \(l\)
 e têm de ser não nulos.

 \item \(A,A_1,B,V_j,L_k,R_k\) têm de descodificar como pontos comprimidos
 canónicos sobre a curva. Uma codificação não canónica ou que não representa
 um ponto provoca rejeição da prova.

 \item Os pontos fornecidos pela prova são multiplicados pelo cofactor 8 antes
 da multiplicação multiescalar. Não há, contudo, uma regra separada que rejeite
 toda ocorrência da identidade: uma identidade honesta é extremamente
 improvável nos pontos aleatórios, mas o verificador decide-a pela estrutura e
 pela equação completa. Os geradores derivados são explicitamente não nulos.

 \item A prova tem de ser não vazia, os vectores
 \(\boldsymbol L,\boldsymbol R\) têm o mesmo comprimento exacto e esse
 comprimento tem de concordar com o número de compromissos e com o limite de
 16. A transacção RingCT deste tipo tem exactamente uma prova canónica.

 \item A equação de lote, a correspondência entre prova e
 \texttt{outPk}, a Equação~\eqref{eq:commitment-balance} e as CLSAG são testes
 distintos; todos têm de passar.
\end{itemize}


Chegámos assim à afirmação prometida: para cada compromisso ligado ao traslado
de \(\Pi_+\), o provador conhece uma abertura cujo montante está em
\([0,2^{64}-1]\). Continuam públicos o número de saídas, os compromissos, o
formato e tamanho da prova e a taxa da transacção. O próximo capítulo juntará
esta propriedade às condições de autorização das entradas, de reconhecimento
das saídas, de equilíbrio e de ausência de dupla despesa que compõem a
validade de uma transacção.
